\documentclass[12pt,a4paper]{article}

\usepackage{iftex}
\ifPDFTeX
  \usepackage[utf8]{inputenc}
  \usepackage[T5]{fontenc}
  \usepackage[vietnamese]{babel}
  \usepackage{lmodern}
\else
  % phương án dự phòng khi biên dịch bằng XeLaTeX/LuaLaTeX
  \usepackage{fontspec}
  \setmainfont[BoldFont=lmroman10-bold.otf,ItalicFont=lmroman10-italic.otf,BoldItalicFont=lmroman10-bolditalic.otf]{lmroman10-regular.otf}\setsansfont{lmsans10-regular.otf}\setmonofont{lmmono10-regular.otf}
  \renewcommand{\contentsname}{Mục lục}\renewcommand{\figurename}{Hình}\renewcommand{\tablename}{Bảng}
\fi
\usepackage[margin=2.2cm]{geometry}
\usepackage{amssymb}
\usepackage{graphicx}
\usepackage{float}
\PassOptionsToPackage{table}{xcolor}
\usepackage{tikz}
\usetikzlibrary{calc,shapes.geometric,arrows.meta,positioning}
\usepackage{booktabs}
\usepackage{longtable}
\usepackage{tabularx}
\usepackage{array}
\usepackage[table]{xcolor}
\usepackage{listings}
\usepackage{enumitem}
\usepackage{fancyhdr}
\usepackage{titlesec}
\usepackage[hidelinks]{hyperref}

\definecolor{bkblue}{RGB}{3,43,145}
\definecolor{codebg}{RGB}{246,248,250}
\newcommand{\todo}[1]{\textcolor{red}{\textbf{[#1]}}}
\newcolumntype{L}[1]{>{\raggedright\arraybackslash}p{#1}}

\lstset{
  basicstyle=\ttfamily\footnotesize, backgroundcolor=\color{codebg},
  frame=single, rulecolor=\color{black!20},
  breaklines=true, showstringspaces=false, columns=fullflexible
}

\titleformat{\section}{\Large\bfseries\color{bkblue}}{\thesection.}{0.6em}{}
\titleformat{\subsection}{\large\bfseries}{\thesubsection.}{0.6em}{}

\pagestyle{fancy}
\fancyhf{}
\lhead{Kiểm tra phần mềm -- Nhóm HEHE2027}
\rhead{Bài tập lớn 3}
\cfoot{\thepage}
\setlength{\headheight}{15pt}

\renewcommand{\arraystretch}{1.2}
\sloppy
\setlength{\emergencystretch}{3em}

\begin{document}

\begin{titlepage}
\begin{tikzpicture}[overlay,remember picture]
\draw [line width=2.25pt]
    ($ (current page.north west) + (2.0cm,-1.5cm) $)
    rectangle
    ($ (current page.south east) + (-1.5cm,2cm) $);
\draw [line width=0.5pt]
    ($ (current page.north west) + (2.2cm,-1.7cm) $)
    rectangle
    ($ (current page.south east) + (-1.7cm,2.2cm) $);
\end{tikzpicture}
\begin{center}
\textbf{ĐẠI HỌC QUỐC GIA THÀNH PHỐ HỒ CHÍ MINH} \\
\textbf{TRƯỜNG ĐẠI HỌC BÁCH KHOA}\\
\textbf{KHOA KHOA HỌC VÀ KỸ THUẬT MÁY TÍNH}

\vspace{1cm}
\includegraphics[width=0.32\textwidth]{img/logo/logo.png}
\vspace{1cm}

\begin{tabular}{c}
    \textbf{\Large KIỂM TRA PHẦN MỀM (SOFTWARE TESTING)}\\[0.4cm]
    \textbf{\large BÀI TẬP LỚN 3: DATA-DRIVEN AUTOMATION TESTING}\\
    \textbf{\large VÀ NON-FUNCTIONAL TESTING}\\[0.2cm]
    \large Kiểm thử hệ thống quản lý học tập \textit{Moodle LMS}\\[0.4cm]
    \textbf{Học kỳ 1, năm học 2026--2027 -- Lớp L01 -- Nhóm HEHE2027}\\
    \hline
\end{tabular}

\vspace{1cm}
\begin{tabular}{p{3cm} p{6cm} l}
    \textbf{GVHD:} & Trần Thị Ngọc Trâm & \\[0.2cm]
    \textbf{Sinh viên:} & Nguyễn Đức Phương & 2312749 \\
                        & Lê Sĩ Nhật Quang & 2312791 \\
                        & Huỳnh Đức Tài & 2212971 \\
                        & Đào Xuân Sơn & 2312954 \\
                        & Sơn Văn Phát & 2212535 \\
\end{tabular}

\vfill
\fontsize{14pt}{0pt}\selectfont TP. HỒ CHÍ MINH, THÁNG 11/2026
\end{center}
\end{titlepage}

\tableofcontents
\newpage

\section{Giới thiệu}
Bài tập lớn 3 tự động hóa lại các test case của Bài tập lớn 2 bằng \textbf{kiểm thử hướng dữ liệu}
(\textit{data-driven testing}) viết bằng \textbf{Python 3}, \textbf{Selenium WebDriver} và \textbf{unittest},
đồng thời thực hiện \textbf{kiểm thử phi chức năng} (\textit{non-functional testing}) trên Moodle LMS.

\begin{itemize}[noitemsep]
  \item \textbf{Level 1:} các test case chỉ khác nhau về dữ liệu được gộp thành một script; dữ liệu đầu vào \emph{và}
        kết quả mong đợi (\textit{Expected Result}) nằm trong file \texttt{.csv}.
  \item \textbf{Level 2:} ngoài dữ liệu, \textbf{test item} (URL của site, locator của ô nhập, nút bấm\ldots) cũng được đưa ra
        file cấu hình/dữ liệu, script không còn chứa hằng số về giao diện.
  \item \textbf{Non-functional:} mỗi thành viên kiểm thử ít nhất một yêu cầu phi chức năng; báo cáo nêu loại kiểm thử, cách tiếp cận và công cụ.
\end{itemize}

\textbf{Thang điểm của đề:} Level 1 -- 5 điểm; Level 2 -- 2 điểm; Non-functional -- 2 điểm; báo cáo -- 1 điểm.

\subsection*{Cấu trúc file nộp (BK LMS)}
Bài nộp gồm một file báo cáo PDF và một file nén \texttt{.zip} cùng tên; file nén chứa năm thư mục, mỗi thư mục là một feature của một thành viên, trong đó file Level 1 và Level 2 được tách riêng:
\begin{lstlisting}
ASSIGNMENT_3__HEHE2027__L01.pdf          (bao cao cua ca nhom)
ASSIGNMENT_3__HEHE2027__L01.zip
  |-- F001/                              (Add a new user - Phuong)
  |     |-- level1/
  |     |     l1_common.py, test_create_user_*.py (3 file), data_L1_*.csv (3 file)
  |     |-- level2/
  |     |     l2_common.py, test_create_user_*.py (3 file), data_L2_*.csv (3 file)
  |     |     config.csv, locators.csv, form_fields.csv
  |     |-- nonfunctional/
  |           perf_create_user.py
  |-- F002/                              (Add a new course - Quang)
  |     |-- level1/   |-- level2/   |-- nonfunctional/
  |-- F003/                              (Tai)
  |     |-- level1/   |-- level2/   |-- nonfunctional/
  |-- F004/                              (Add new category - Son)
  |     |-- level1/   |-- level2/   |-- nonfunctional/
  |-- F005/                              (Add a Choice - Phat)
        |-- level1/   (test_choice.py, engine.py, settings.py, choice_data.csv)
        |-- level2/   (test_choice.py, engine.py, choice_data.csv,
        |             environment.csv, ui_elements.csv)
        |-- nonfunctional/ (test_compatibility.py)
        run_all.py, CHAY_TAT_CA.cmd, README.md
\end{lstlisting}
Chỉ nộp mã Python và file dữ liệu, \textbf{không} nộp file Katalon. Các script dùng hàm chung trong chính thư mục feature: Feature 001 dùng \texttt{l1\_common.py}/\texttt{l2\_common.py}; Feature 005 dùng \texttt{engine.py}. Cách chạy cụ thể trình bày ở từng feature. Với Feature 005, giữ các thư mục \texttt{level1}, \texttt{level2}, \texttt{nonfunctional} cạnh nhau để script tương thích gọi lại mã Level 2.

\section{Phân công công việc}
{\small
\begin{longtable}{|c|L{3.1cm}|L{3.4cm}|L{4.3cm}|c|}
\hline
\rowcolor{gray!15}
\textbf{STT} & \textbf{Thành viên} & \textbf{Feature (BTL2)} & \textbf{Công việc BTL3} & \textbf{Đóng góp} \\ \hline
\endhead
1 & Nguyễn Đức Phương \newline 2312749 & 001 -- Add a new user & Level 1, Level 2 cho feature 001; kiểm thử phi chức năng (mục \ref{sec:nf}) & 100\%  \\ \hline
2 & Lê Sĩ Nhật Quang \newline 2312791 & 002 -- Add a new course & Level 1, Level 2 cho feature 002; phi chức năng & \todo{\%} \\ \hline
3 & Huỳnh Đức Tài \newline 2212971 & \todo{chưa có feature} & \todo{điền} & \todo{\%} \\ \hline
4 & Đào Xuân Sơn \newline 2312954 & 004 -- Add new category & Level 1, Level 2 cho feature 004; phi chức năng & \todo{\%} \\ \hline
5 & Sơn Văn Phát \newline 2212535 & 005 -- Add a Choice & Level 1, Level 2 cho feature 005; phi chức năng & 100\% \\ \hline
\end{longtable}
}

\section{Môi trường và công cụ}
\begin{itemize}[noitemsep]
  \item Moodle 5.1 chạy bằng Docker, \texttt{http://127.0.0.1:9501}, tài khoản \texttt{admin}/\texttt{admin}. Không đổi mật khẩu các user mặc định.
  \item Python 3.x, Selenium 4 (Selenium Manager tự tải driver), module chuẩn \texttt{unittest}, \texttt{csv}; Chrome, Firefox. Feature 005 chạy trên Chrome và Edge.
  \item Katalon Recorder chỉ dùng để thu thao tác và tham khảo locator; sản phẩm nộp là script Python.
  \item Ngôn ngữ giao diện Moodle: English.
\end{itemize}

\section{Feature 001: Add a new user -- kiểm thử hướng dữ liệu}

\subsection{Phân tích gộp 50 test case thành các script}
Quan sát: 45 test case đầu của BTL2 (ECP, BVA, decision table) có cùng một quy trình --
\emph{mở form, điền các trường, bấm Create user, kiểm tra kết quả} -- chỉ khác dữ liệu và kết quả mong đợi.
Các use-case còn lại khác về luồng thao tác nên tách riêng. Kết quả gộp:

\begin{table}[H]
\centering\small
\begin{tabular}{|L{4.4cm}|L{2.7cm}|L{4.6cm}|c|}
\hline
\rowcolor{gray!15}\textbf{Script (Level 1)} & \textbf{Test case gộp} & \textbf{Luồng} & \textbf{Số dòng dữ liệu} \\ \hline
\texttt{test\_create\_user\_form} & TC-001-001 đến 045 & Điền form, Create user, so với Expected Result (thành công hoặc lỗi tại một trường). & 45 \\ \hline
\texttt{test\_create\_user\_retry} & TC-001-047, 048 & Gửi dữ liệu lần 1 (lỗi), sửa trường lỗi, gửi lần 2 (thành công). & 2 \\ \hline
\texttt{test\_create\_user\_nav} & TC-001-046, 049, 050 & Luồng điều hướng: từ danh sách user vào form, Cancel, hoặc tick Generate password (khóa ô New password). & 3 \\ \hline
\end{tabular}
\caption{Gộp 50 test case thành 3 script}
\end{table}
Tính ra, 50 test case thủ công được thay bằng 3 script và 3 file dữ liệu; thêm một test case mới chỉ là thêm một dòng vào file \texttt{.csv}.

\subsection{Cấu trúc file dữ liệu Level 1}
File \texttt{data\_L1\_create\_user.csv} (một dòng = một test case):

\begin{table}[H]
\centering\footnotesize
\begin{tabular}{|L{5.0cm}|L{9.6cm}|}
\hline
\rowcolor{gray!15}\textbf{Cột} & \textbf{Ý nghĩa} \\ \hline
\texttt{tc\_id} & Mã test case (ví dụ \texttt{TC-001-003}), dùng làm tên \texttt{subTest}. \\ \hline
\texttt{username, password, firstname, lastname, email, city} & Dữ liệu nhập. Chuỗi dài viết bằng ký hiệu \texttt{<a*101>} (101 ký tự \texttt{a}), script mở rộng khi chạy. \\ \hline
\texttt{generate\_password, force\_change, suspended} & \texttt{0/1}: trạng thái ba checkbox. \\ \hline
\texttt{expected\_type} & \texttt{success} (có thông báo ``Changes saved''; thường là về trang danh sách, riêng Generate password thì ở lại form) hoặc \texttt{error} (ở lại form, hiện lỗi). \\ \hline
\texttt{error\_field} & Trường báo lỗi (\texttt{username}, \texttt{newpassword}, \texttt{firstname}, \texttt{lastname}, \texttt{email}\ldots); rỗng nếu thành công. \\ \hline
\texttt{expected\_text} & Nội dung thông báo mong đợi (ví dụ \texttt{This username already exists, choose another}). \\ \hline
\end{tabular}
\caption{Cột của file dữ liệu Level 1}
\end{table}

Ví dụ một số dòng (rút gọn):
\begin{lstlisting}
tc_id,username,password,firstname,lastname,email,city,generate_password,force_change,suspended,expected_type,error_field,expected_text
TC-001-001,hehe001,Test@1234,An,Nguyen,hehe001@example.com,HCM,0,0,0,success,,Changes saved
TC-001-002,,Test@1234,An,Nguyen,hehe002@example.com,,0,0,0,error,username,Required
TC-001-004,HeheUPPER,Test@1234,An,Nguyen,hehe004@example.com,,0,0,0,error,username,Only lowercase letters allowed
TC-001-006,admin,Test@1234,An,Nguyen,hehe006@example.com,,0,0,0,error,username,This username already exists
TC-001-014,hehe014,Test@1234,An,Nguyen,khong-co-a-cong,,0,0,0,error,email,Invalid email address
TC-001-029,hehe029,Test@1234,<a*100>,Nguyen,hehe029@example.com,,0,0,0,success,,Changes saved
\end{lstlisting}

\subsection{Thiết kế script Level 1}
Mã đầy đủ nằm trong thư mục \texttt{F001/level1/} (hàm dùng chung ở \texttt{l1\_common.py}, ba script \texttt{test\_create\_user\_*.py}); đoạn dưới đây là bản rút gọn. Script dùng \texttt{unittest.TestCase}; mỗi dòng CSV chạy trong một \texttt{subTest}, nên một dòng sai không làm dừng các dòng còn lại và báo cáo ghi rõ test case nào fail.

\begin{lstlisting}[language=Python]
import csv, os, re, unittest
from selenium import webdriver
from selenium.webdriver.common.by import By
from selenium.webdriver.support.ui import WebDriverWait
from selenium.webdriver.support import expected_conditions as EC

BASE = "http://127.0.0.1:9501"
DATA = os.path.join(os.path.dirname(__file__), "data_L1_create_user.csv")

def expand(v):                      # "<a*101>" -> "a" * 101
    return re.sub(r"<(.)\*(\d+)>", lambda m: m.group(1) * int(m.group(2)), v)

def read_rows(path):
    with open(path, newline="", encoding="utf-8") as f:
        return [{k: expand(v) for k, v in r.items()} for r in csv.DictReader(f)]

class CreateUserForm(unittest.TestCase):
    @classmethod
    def setUpClass(cls):
        cls.d = webdriver.Chrome()
        cls.wait = WebDriverWait(cls.d, 10)
        cls.d.get(BASE + "/login/index.php")           # buoc chuan bi, khong phai doi tuong test
        cls.d.find_element(By.ID, "username").send_keys("admin")
        cls.d.find_element(By.ID, "password").send_keys("admin")
        cls.d.find_element(By.ID, "loginbtn").click()

    @classmethod
    def tearDownClass(cls):
        cls.d.quit()

    def set_checkbox(self, elem_id, want):
        box = self.d.find_element(By.ID, elem_id)
        if box.is_selected() != (want == "1"):
            box.click()

    def fill_form(self, r):
        d = self.d
        d.get(BASE + "/user/editadvanced.php?id=-1")
        for fid in ("username", "firstname", "lastname", "email", "city"):
            e = d.find_element(By.ID, "id_" + fid); e.clear(); e.send_keys(r[fid])
        self.set_checkbox("id_createpassword", r["generate_password"])
        if r["generate_password"] != "1":              # o password bi an, gan bang JS
            d.execute_script("arguments[0].value = arguments[1];",
                             d.find_element(By.ID, "id_newpassword"), r["password"])
        self.set_checkbox("id_preference_auth_forcepasswordchange", r["force_change"])
        self.set_checkbox("id_suspended", r["suspended"])
        d.find_element(By.ID, "id_submitbutton").click()

    def test_create_user(self):
        for r in read_rows(DATA):
            with self.subTest(tc=r["tc_id"]):
                self.fill_form(r)
                if r["expected_type"] == "success":
                    self.wait.until(lambda x: "/admin/user.php" in x.current_url
                                    or "Changes saved" in x.page_source)
                    self.assertIn(r["expected_text"], self.d.page_source)
                else:
                    err = self.wait.until(EC.visibility_of_element_located(
                        (By.ID, "id_error_" + r["error_field"])))
                    self.assertIn(r["expected_text"], err.text)

if __name__ == "__main__":
    unittest.main(verbosity=2)
\end{lstlisting}

\noindent\textbf{Lưu ý thiết kế:}
\begin{itemize}[noitemsep]
  \item \textbf{Dữ liệu độc nhất:} username/email trong CSV là duy nhất cho mỗi dòng và mỗi lần chạy phải xóa user đã tạo (hoặc đổi tiền tố); nếu không, lần chạy sau sẽ báo trùng.
  \item \textbf{Giới hạn \texttt{maxlength}:} ô First name/Last name/Email/City có \texttt{maxlength}, nên chuỗi 101 ký tự bị trình duyệt cắt còn 100; Expected Result của các dòng này phải khớp hành vi thực tế (ví dụ \texttt{success}).
  \item \textbf{Hộp thoại ``Reload site?''} do Moodle bật khi form đang sửa dở: script phải tắt cảnh báo này (hoặc dùng \texttt{get} sau khi đã submit) trước khi mở trang mới.
  \item Ô New password bị ẩn nên gán bằng JavaScript (không bị trình duyệt tự cắt); script tự cắt giá trị theo thuộc tính \texttt{maxlength} (128) để mô phỏng người dùng gõ, nên dòng 129 ký tự được kỳ vọng là \texttt{success} với mật khẩu 128 ký tự.
\end{itemize}

\subsection{Level 2: đưa test item ra ngoài script}
Ở Level 2, script không còn chứa URL hay locator. Ba file cấu hình bổ sung:

\begin{lstlisting}
config.csv        key,value
                  base_url,http://127.0.0.1:9501
                  login_path,/login/index.php
                  form_path,/user/editadvanced.php?id=-1
                  admin_user,admin
                  admin_pass,admin

locators.csv      name,by,value
                  username_input,id,id_username
                  submit_button,id,id_submitbutton
                  error_username,id,id_error_username
                  ...

form_fields.csv   data_column,locator,kind
                  username,username_input,text
                  password,password_input,password_js
                  suspended,suspended_checkbox,checkbox
\end{lstlisting}
Script Level 2 đọc \texttt{form\_fields.csv} theo thứ tự và với mỗi dòng gọi hành động tương ứng với \texttt{kind} (\texttt{text}, \texttt{password\_js}, \texttt{checkbox}); locator được tra trong \texttt{locators.csv} bằng \texttt{by} (\texttt{id}, \texttt{name}, \texttt{xpath}, \texttt{css}) ánh xạ sang \texttt{selenium.webdriver.common.by.By}. Cột \texttt{error\_field} của file dữ liệu cũng là một \textbf{tên locator} (ví dụ \texttt{error\_username}) chứ không còn là chuỗi id cố định.

\begin{lstlisting}[language=Python]
LOCATORS = {r["name"]: (BY[r["by"]], r["value"]) for r in read_csv("locators.csv")}
FIELDS   = read_csv("form_fields.csv")
CFG      = {r["key"]: r["value"] for r in read_csv("config.csv")}

def fill_form(self, row):
    self.d.get(CFG["base_url"] + CFG["form_path"])
    for f in FIELDS:
        by, val = LOCATORS[f["locator"]]
        el = self.d.find_element(by, val)
        if f["kind"] == "text":
            el.clear(); el.send_keys(row[f["data_column"]])
        elif f["kind"] == "password_js":
            self.d.execute_script("arguments[0].value = arguments[1];", el, row[f["data_column"]])
        elif f["kind"] == "checkbox":
            if el.is_selected() != (row[f["data_column"]] == "1"): el.click()
\end{lstlisting}
Khi giao diện đổi (ví dụ id của ô nhập thay đổi, hoặc nhóm chạy ở site khác) chỉ cần sửa \texttt{locators.csv} hoặc \texttt{config.csv}, không sửa script. Thêm trường mới vào form chỉ cần thêm một dòng vào \texttt{form\_fields.csv} và một cột vào file dữ liệu.

\begin{table}[H]
\centering\small
\begin{tabular}{|L{3.4cm}|L{5.4cm}|L{5.4cm}|}
\hline
\rowcolor{gray!15} & \textbf{Level 1} & \textbf{Level 2} \\ \hline
Dữ liệu và Expected Result & Trong CSV & Trong CSV \\ \hline
URL site & Hằng số trong script & \texttt{config.csv} \\ \hline
Locator ô nhập, nút bấm & Hằng số trong script & \texttt{locators.csv} \\ \hline
Danh sách trường của form & Viết cứng trong \texttt{fill\_form} & \texttt{form\_fields.csv} \\ \hline
\end{tabular}
\caption{So sánh Level 1 và Level 2}
\end{table}

\subsection{Cách chạy}
\begin{lstlisting}
pip install selenium
cd level1
python -m unittest -v test_create_user_form test_create_user_retry test_create_user_nav
\end{lstlisting}
Biến môi trường tùy chọn: \texttt{RUN\_ID} (đặt mã lần chạy), \texttt{KEEP\_DATA=1} (không xóa user sau khi chạy). Cuối mỗi script có bước dọn dữ liệu: xóa các user có username/email chứa \texttt{tc001\_} hoặc \texttt{a@b.co} qua trang \texttt{/admin/user.php}; Level 2 chạy tương tự trong thư mục \texttt{F001/level2/}.

\subsection{Kết quả chạy}
Chạy ngày 21/11/2026 trên Chrome for Testing 158, Moodle 5.1 demo (\texttt{http://127.0.0.1:9501}), lệnh
\texttt{python -m unittest -v test\_create\_user\_form test\_create\_user\_retry test\_create\_user\_nav} trong thư mục \texttt{level1/}.

\begin{table}[H]
\centering\small
\begin{tabular}{|L{5.0cm}|L{3.6cm}|c|c|}
\hline
\rowcolor{gray!15}\textbf{Script} & \textbf{Test case} & \textbf{Số dòng} & \textbf{Kết quả} \\ \hline
\texttt{test\_create\_user\_form} & TC-001-001 đến 045 & 45 & 45 pass \\ \hline
\texttt{test\_create\_user\_retry} & TC-001-047, 048 & 2 & 2 pass \\ \hline
\texttt{test\_create\_user\_nav} & TC-001-046, 049, 050 & 3 & 3 pass \\ \hline
\end{tabular}
\caption{Kết quả chạy Level 1: 50/50 test case pass, tổng thời gian khoảng 367 giây}
\end{table}

\textbf{Các điều chỉnh rút ra khi chạy lần đầu:}
\begin{itemize}[noitemsep]
  \item TC-001-004 và TC-001-005: thông báo thực tế của Moodle là \texttt{Only lowercase letters allowed} và \texttt{The username can only contain alphanumeric lowercase characters\ldots}; Expected Result trong file dữ liệu được sửa theo đó.
  \item TC-001-041 (Generate password): Moodle ở lại trang form và hiện thông báo \texttt{Changes saved} thay vì chuyển sang danh sách; điều kiện thành công của script được sửa thành ``chuyển sang danh sách hoặc có \texttt{Changes saved}''.
  \item Nút Create user nằm cuối trang dài và bị che bởi footer dính nên script cuộn phần tử vào giữa màn hình trước khi bấm; đăng nhập được thử lại và kiểm tra giá trị ô nhập để tránh gõ thiếu khi trang chưa tải xong.
\end{itemize}
\textbf{Level 2:} chạy cùng lệnh trong thư mục \texttt{level2/}, với URL, locator và danh sách trường đọc từ \texttt{config.csv}, \texttt{locators.csv}, \texttt{form\_fields.csv}.

\begin{table}[H]
\centering\small
\begin{tabular}{|L{5.0cm}|L{3.6cm}|c|c|}
\hline
\rowcolor{gray!15}\textbf{Script} & \textbf{Test case} & \textbf{Số dòng} & \textbf{Kết quả} \\ \hline
\texttt{test\_create\_user\_form} & TC-001-001 đến 045 & 45 & 45 pass \\ \hline
\texttt{test\_create\_user\_retry} & TC-001-047, 048 & 2 & 2 pass \\ \hline
\texttt{test\_create\_user\_nav} & TC-001-046, 049, 050 & 3 & 3 pass \\ \hline
\end{tabular}
\caption{Kết quả chạy Level 2: 50/50 test case pass, tổng thời gian khoảng 362 giây}
\end{table}

\section{Feature 002}
\section{Feature 003}
\section{Feature 004}
\section{Feature 005: Add a Choice -- kiểm thử hướng dữ liệu}

\subsection{Phân tích gộp 24 test case thành script}
Các test case của Feature 005 ở BTL2 cùng thao tác trên form Add a Choice, nhưng khác tên hoạt động,
phương án trả lời và các thiết lập. Vì vậy, hai mức kiểm thử dùng chung một luồng xử lý trong
\texttt{engine.py}; dữ liệu và kết quả mong đợi của 24 ca nằm trong \texttt{choice\_data.csv}.
Ba nhóm luồng được phân biệt bằng giá trị \texttt{outcome}:

\begin{table}[H]
\centering\small
\begin{tabular}{|L{2.7cm}|L{4.0cm}|L{5.0cm}|c|}
\hline
\rowcolor{gray!15}\textbf{Nhóm} & \textbf{Test case gộp} & \textbf{Luồng} & \textbf{Số dòng} \\ \hline
Tạo hợp lệ & TC-005-001, 003--006, 008--019, 021--023 & Nhập dữ liệu, lưu, kiểm tên và mở lại Settings để đối chiếu cấu hình. & 20 \\ \hline
Báo lỗi rồi sửa & TC-005-002, 007, 020 & Gửi dữ liệu sai, kiểm thông báo lỗi, sửa trường lỗi rồi lưu lại. & 3 \\ \hline
Hủy tạo & TC-005-024 & Nhập dữ liệu, bấm Cancel, kiểm danh sách hoạt động không thay đổi. & 1 \\ \hline
\end{tabular}
\caption{Gộp 24 test case của Add a Choice theo luồng xử lý}
\end{table}

Khác với cách tách ba script của Feature 001, Feature 005 sử dụng một script
\texttt{test\_choice.py} cho cả ba nhóm. Khi chạy, mỗi dòng CSV được tạo thành một test
\texttt{unittest} riêng, có tên tương ứng với mã ca. Một ca lỗi không làm dừng các ca còn lại.

\subsection{Cấu trúc file dữ liệu Level 1}
File \texttt{F005/level1/choice\_data.csv} có 24 dòng dữ liệu, mã hóa UTF-8 BOM để mở được tiếng Việt
trong Excel. Các nhóm giá trị được lưu dưới dạng chuỗi JSON trong ô CSV, do số phương án và trường
thiết lập không giống nhau giữa các ca.

\begin{table}[H]
\centering\small
\begin{tabular}{|L{3.6cm}|L{10.5cm}|}
\hline
\rowcolor{gray!15}\textbf{Cột} & \textbf{Ý nghĩa} \\ \hline
\texttt{case\_id} & Mã ca, từ TC-005-001 đến TC-005-024; dùng trong tên test và file kết quả. \\ \hline
\texttt{case\_name} & Nội dung kiểm thử, ví dụ tên rỗng, giới hạn lượt chọn hoặc Cancel. \\ \hline
\texttt{fields\_json} & Dữ liệu nhập lần đầu: tên, các phương án, mô tả, giới hạn và thời gian mở/đóng. \\ \hline
\texttt{repair\_json} & Dữ liệu sửa sau khi kiểm lỗi; rỗng đối với ca không có bước sửa. \\ \hline
\texttt{expected\_json} & Kết quả mong đợi: loại luồng, tên, số phương án, thông báo lỗi và các giá trị cần đọc lại sau khi lưu. \\ \hline
\end{tabular}
\caption{Cột của file dữ liệu Level 1 -- Feature 005}
\end{table}

Trong \texttt{expected\_json}, \texttt{outcome} nhận một trong ba giá trị:
\texttt{created}, \texttt{error\_then\_created}, \texttt{cancelled}.
Ví dụ dưới đây biểu diễn dữ liệu của TC-005-002 sau khi đọc và giải mã các ô CSV, đã lược bớt
các thiết lập giữ mặc định:

\begin{lstlisting}[language=Python]
case_id = "TC-005-002"
fields = {"name": ""}  # Other input fields omitted.
repair = {"name": "F005-CHOICE-002"}
expected = {
    "outcome": "error_then_created",
    "error_field": "name",
    "error_text": "You must supply a value here.",
    "title": "F005-CHOICE-002",
    "option_count": 2,
    "saved_fields": {"name": "F005-CHOICE-002", "...": "..."}
}
\end{lstlisting}
Hai chuỗi tên ở TC-005-005 và TC-005-006 được lưu đầy đủ trong CSV, lần lượt có 1.332 và 1.333
ký tự. Các thiết lập khác gồm hiển thị dọc, cho phép thay đổi câu trả lời, chọn nhiều phương án,
giới hạn 0/1/2, giờ mở/đóng, cột người chưa trả lời và công bố kết quả ẩn danh.

\subsection{Thiết kế script Level 1}
Mã nằm trong \texttt{F005/level1/}: \texttt{test\_choice.py} là điểm chạy,
\texttt{engine.py} chứa phần xử lý dùng chung, còn \texttt{settings.py} chứa URL và locator.
Script đọc CSV bằng \texttt{csv.DictReader}, giải mã các nhóm giá trị bằng \texttt{json.loads},
rồi chạy cùng hàm \texttt{scenario} cho từng dòng.

Đoạn dưới rút gọn cách kiểm tra ba nhóm luồng; các hàm thao tác trang được cài đặt trong
\texttt{ChoicePage}:

\begin{lstlisting}[language=Python]
def scenario(self, row):
    page, expected = self.page, row["expected"]
    page.edit_mode(True)
    before = page.activity_ids()
    page.go("create_url")
    page.fields(row["fields"])

    if expected["outcome"] == "cancelled":
        page.click("cancel")
        page.wait.until(lambda d: "/course/view.php" in d.current_url)
        self.assertEqual(before, page.activity_ids())
        return

    page.click("save")
    if expected["outcome"] == "error_then_created":
        error = page.element("error", field=expected["error_field"])
        self.assertEqual(expected["error_text"], error.text.strip())
        page.fields(row["repair"])
        page.click("save")

    cmid = page.wait_created()
    self.assertEqual(expected["title"],
                     page.element("activity_heading").text.strip())
    page.go("settings_url", cmid=cmid)
    for key, value in expected["saved_fields"].items():
        self.assertEqual(value, page.read_field(key))
\end{lstlisting}

\noindent\textbf{Lưu ý thiết kế:}
\begin{itemize}[noitemsep]
  \item \textbf{Đọc lại dữ liệu:} sau khi lưu, mở Settings của hoạt động vừa tạo để đối chiếu
        giá trị thực với \texttt{saved\_fields}; kiểm thêm số phương án đã lưu. Các ca 016, 017 và 024
        đều có phép kiểm tra kết quả, thay cho phần còn thiếu trong bản ghi Katalon.
  \item \textbf{Chờ đúng trạng thái:} dùng \texttt{WebDriverWait} để chờ trang đổi sau khi bật/tắt Edit mode,
        chờ lỗi hiển thị và chờ hoạt động bị xóa. Không dùng ID YUI sinh theo lần tải trang.
  \item \textbf{Mô tả:} nhập văn bản trong iframe TinyMCE. TC-005-010 còn kiểm mô tả xuất hiện
        tại đúng hoạt động trên trang khóa học.
  \item \textbf{Dọn dữ liệu:} lấy ID hoạt động từ URL sau khi tạo. Trong \texttt{tearDown}, chỉ xóa
        hoạt động của ca vừa chạy, kiểm đã xóa xong rồi tắt Edit mode. Giữ nguyên khóa học và dữ liệu có sẵn.
        Lỗi dọn dữ liệu được ghi nhận là \texttt{ERROR}, không tính là pass.
\end{itemize}

\subsection{Level 2: đưa test item ra ngoài script}
Level 2 tiếp tục dùng \texttt{choice\_data.csv}, đồng thời chuyển cấu hình từ
\texttt{settings.py} sang hai file CSV:

\begin{table}[H]
\centering\small
\begin{tabular}{|L{4.2cm}|L{9.9cm}|}
\hline
\rowcolor{gray!15}\textbf{File} & \textbf{Nội dung} \\ \hline
\texttt{environment.csv} & Các cặp \texttt{key,value}: địa chỉ Moodle, ID và tên khóa học,
thời gian chờ, URL đăng nhập, khóa học, form tạo Choice và Settings. \\ \hline
\texttt{ui\_elements.csv} & Các cột \texttt{key,by,selector,kind}: tên phần tử, loại locator,
selector và kiểu thao tác như nhập văn bản, chọn giá trị, checkbox, ngày giờ hoặc editor. \\ \hline
\end{tabular}
\caption{File cấu hình bổ sung của Level 2 -- Feature 005}
\end{table}

Ví dụ dữ liệu cấu hình, rút gọn:
\begin{lstlisting}
environment.csv   key,value
                  base_url,http://127.0.0.1:9501
                  course_id,2
                  course_name,F005
                  course_url,{base_url}/course/view.php?id={course_id}

ui_elements.csv   key,by,selector,kind
                  name,css,#id_name,text
                  option,css,#id_option_{index},text
                  allowmultiple,css,#id_allowmultiple,select
                  save,css,#id_submitbutton,text
\end{lstlisting}

\begin{lstlisting}[language=Python]
CONFIG = {row["key"]: row["value"]
          for row in read_csv(ROOT / "environment.csv")}
UI = {row["key"]: row
      for row in read_csv(ROOT / "ui_elements.csv")}
run(CONFIG, UI, ROOT)
\end{lstlisting}

Hàm thao tác tra locator theo tên phần tử trong \texttt{UI}; các mẫu
\texttt{\{index\}} và \texttt{\{cmid\}} được thay bằng chỉ số phương án hoặc ID hoạt động thực tế.
Khi thay địa chỉ Moodle hoặc selector của phần tử đã hỗ trợ, chỉ cần sửa CSV tương ứng.

\begin{table}[H]
\centering\small
\begin{tabular}{|L{4.0cm}|L{4.8cm}|L{5.0cm}|}
\hline
\rowcolor{gray!15} & \textbf{Level 1} & \textbf{Level 2} \\ \hline
Dữ liệu và Expected Result & \texttt{choice\_data.csv} & \texttt{choice\_data.csv} \\ \hline
URL và cấu hình khóa học & \texttt{settings.py} & \texttt{environment.csv} \\ \hline
Locator và kiểu phần tử & \texttt{settings.py} & \texttt{ui\_elements.csv} \\ \hline
Luồng và phép kiểm tra & \texttt{engine.py} & \texttt{engine.py} \\ \hline
\end{tabular}
\caption{So sánh hai mức triển khai của Feature 005}
\end{table}

\subsection{Cách chạy}
\textbf{Tiền điều kiện:} Moodle đã khởi động, tài khoản quản trị đăng nhập được và có khóa học
tên đầy đủ F005. Trên Moodle mới, tạo khóa một lần tại
\emph{Site administration > Courses > Add a new course}, sau đó lấy ID từ URL
\texttt{/course/view.php?id=N}. Đây là bước chuẩn bị môi trường, không nằm trong 24 ca Add a Choice.
Bộ mã được kiểm chứng trong báo cáo này chưa tự tạo khóa học.

Từ thư mục \texttt{F005}, cài thư viện rồi chạy từng mức:
\begin{lstlisting}
python -m pip install -r level1/requirements.txt
python level1/test_choice.py --course-id 2
python level2/test_choice.py --course-id 2
\end{lstlisting}
Số 2 là ID của khóa học trong lần chạy được ghi nhận; trên máy khác thay bằng ID thực tế.
Với cấu hình mặc định này có thể chạy \texttt{python run\_all.py}, hoặc bấm
\texttt{CHAY\_TAT\_CA.cmd} trên Windows để chạy lần lượt cả hai Level và phần tương thích.
Khi ID khác 2, chạy từng script với \texttt{--course-id N}; phần phi chức năng dùng cùng tham số.

Chạy riêng một ca hoặc không hiện cửa sổ trình duyệt:
\begin{lstlisting}
python level1/test_choice.py --case TC-005-001
python level2/test_choice.py --headless
\end{lstlisting}
Các tham số tùy chọn gồm \texttt{--base-url}, \texttt{--browser}, \texttt{--keep-created}.
Tài khoản mặc định là admin/admin; có thể đặt \texttt{MOODLE\_USERNAME} và
\texttt{MOODLE\_PASSWORD} khi môi trường dùng tài khoản khác.

\subsection{Kết quả chạy}
Chạy ngày 24/11/2026 trên Windows, Chrome 158.0.8127.37, Selenium 4.54.0,
Moodle demo chạy bằng Docker tại \texttt{http://127.0.0.1:9501}, khóa F005 có ID 2.
Số liệu dưới đây lấy từ lần chạy của Sơn Văn Phát, có đối chiếu \texttt{results.csv},
\texttt{run.json} và kết quả hiển thị trên PowerShell.

\begin{table}[H]
\centering\small
\begin{tabular}{|L{4.7cm}|L{3.5cm}|c|c|}
\hline
\rowcolor{gray!15}\textbf{Nhóm luồng} & \textbf{Script} & \textbf{Level 1} & \textbf{Level 2} \\ \hline
Tạo Choice hợp lệ & \texttt{test\_choice.py} & 20 pass & 20 pass \\ \hline
Báo lỗi rồi sửa & \texttt{test\_choice.py} & 3 pass & 3 pass \\ \hline
Cancel & \texttt{test\_choice.py} & 1 pass & 1 pass \\ \hline
\textbf{Tổng} & \textbf{24 dòng dữ liệu} & \textbf{24/24 pass} & \textbf{24/24 pass} \\ \hline
\end{tabular}
\caption{Kết quả kiểm thử hướng dữ liệu của Feature 005}
\end{table}

Thời gian do \texttt{unittest} ghi nhận là khoảng \textbf{238,128 giây} ở Level 1 và
\textbf{196,423 giây} ở Level 2, bao gồm thao tác từng ca và dọn dữ liệu. Đây là thời gian thực thi
bộ test, không dùng làm tiêu chí đánh giá hiệu năng của Moodle.
Mỗi lần chạy có \texttt{results.csv} ghi trạng thái từng ca và \texttt{run.json} ghi cấu hình,
phiên bản trình duyệt và tổng số pass/fail/error. Trong bộ file đưa lên Drive, hai thư mục
bằng chứng được đặt tên \texttt{F005\_Level1\_Chrome} và \texttt{F005\_Level2\_Chrome}.

\textbf{Các điều chỉnh rút ra khi chạy lần đầu:}
\begin{itemize}[noitemsep]
  \item Chờ trang cũ kết thúc chuyển hướng sau khi đổi Edit mode trước khi mở form tạo Choice.
  \item Cuộn iframe TinyMCE vào vùng nhìn thấy trước khi nhập mô tả, tránh thanh điều hướng che vị trí bấm.
  \item Sau khi xác nhận xóa, chờ hoạt động biến mất rồi mới tải lại trang khóa học; tránh báo lỗi
        do trang được tải lại khi Moodle còn đang xử lý xóa.
\end{itemize}

\section{Kiểm thử phi chức năng}\label{sec:nf}

\begin{table}[H]
\centering\footnotesize
\begin{tabular}{|c|L{2.4cm}|L{3.2cm}|L{3.6cm}|L{3.0cm}|}
\hline
\rowcolor{gray!15}\textbf{Thành viên} & \textbf{Loại kiểm thử} & \textbf{Yêu cầu} & \textbf{Cách tiếp cận} & \textbf{Công cụ} \\ \hline
Phương & Hiệu năng (performance) & Gửi form Create user phản hồi trong $\le$ 3 giây (trung bình) với 20 lần lặp. & Script đo thời gian từ lúc bấm Create user đến khi trang danh sách hiển thị, 20 lần; tính trung bình, lớn nhất. & Selenium + \texttt{time.perf\_counter}, \texttt{Navigation Timing API} \\ \hline
Quang & Tương thích (compatibility) & Form Add a new course hoạt động giống nhau trên Chrome và Firefox. & Chạy lại script Level 1 trên hai trình duyệt, so sánh kết quả từng dòng. & Selenium (Chrome, Firefox) \\ \hline
Tài & \todo{điền} & \todo{điền} & \todo{điền} & \todo{điền} \\ \hline
Sơn & Bảo mật cơ bản (security) & Dữ liệu có ký tự đặc biệt (\verb|<script>|, dấu nháy) không làm hỏng trang và không bị thực thi. & Nhập chuỗi XSS/SQL injection vào Name của category, kiểm tra trang kết quả hiển thị thành văn bản, không có hộp thoại. & Selenium, OWASP payload mẫu \\ \hline
Phát & Tương thích (compatibility) & Bốn luồng Add a Choice đạt và cho dữ liệu đọc lại tương đương trên Chrome/Edge. & Chạy 4 ca trên mỗi trình duyệt, so kết quả kiểm tra và các quan sát đã lưu. & Python, Selenium (Chrome, Edge), \texttt{unittest} \\ \hline
\end{tabular}
\caption{Yêu cầu phi chức năng theo thành viên}
\end{table}

\subsection{Chi tiết phần của Phương: hiệu năng của Create user}
\textbf{Loại kiểm thử:} kiểm thử hiệu năng (thời gian phản hồi) ở mức một người dùng.
\textbf{Cách tiếp cận:} dùng script Level 1 với dữ liệu hợp lệ, mỗi lần dùng username khác nhau; ghi thời gian giữa hai thời điểm: ngay trước khi bấm Create user và khi trang \texttt{/admin/user.php} hiển thị xong. Lặp 20 lần, báo cáo trung bình, nhỏ nhất, lớn nhất.
\textbf{Công cụ:} Selenium WebDriver, \texttt{time.perf\_counter}.

\begin{lstlisting}[language=Python]
t0 = time.perf_counter()
driver.find_element(By.ID, "id_submitbutton").click()
wait.until(EC.url_contains("/admin/user.php"))
elapsed = time.perf_counter() - t0
times.append(elapsed)
\end{lstlisting}
\textbf{Tiêu chí đạt:} thời gian trung bình $\le$ 3 giây.

\textbf{Script và cách chạy:} \texttt{F001/nonfunctional/perf\_create\_user.py} (dùng lại hàm của \texttt{F001/level1/l1\_common.py}). Mỗi lần lặp tạo một user hợp lệ có username/email duy nhất, đo thời gian từ lúc bấm Create user đến lúc hiện ``Changes saved'' hoặc trang danh sách, ghi vào \texttt{perf\_results.csv}, in trung bình, nhỏ nhất, lớn nhất và phân vị 95; test fail nếu trung bình vượt ngưỡng. Sau khi chạy script xóa các user đã tạo.
\begin{lstlisting}
cd nonfunctional
python -m unittest -v perf_create_user
\end{lstlisting}

\textbf{Kết quả đo} (20 lần, chạy ngày 21/11/2026, Chrome for Testing 158, Moodle 5.1 chạy bằng Docker trên cùng máy):

\begin{table}[H]
\centering\small
\begin{tabular}{|c|c|c|c|}
\hline
\rowcolor{gray!15}\textbf{Lần} & \textbf{Create user (s)} & \textbf{Lần} & \textbf{Create user (s)} \\ \hline
1 & 3.205 & 11 & 1.581 \\ \hline
2 & 2.009 & 12 & 1.620 \\ \hline
3 & 2.010 & 13 & 1.623 \\ \hline
4 & 2.155 & 14 & 1.635 \\ \hline
5 & 2.035 & 15 & 1.745 \\ \hline
6 & 2.000 & 16 & 1.680 \\ \hline
7 & 2.061 & 17 & 1.590 \\ \hline
8 & 1.971 & 18 & 1.603 \\ \hline
9 & 1.677 & 19 & 1.649 \\ \hline
10 & 1.661 & 20 & 1.641 \\ \hline
\end{tabular}
\caption{Thời gian từ lúc bấm Create user đến lúc Moodle hiển thị kết quả}
\end{table}

\begin{table}[H]
\centering\small
\begin{tabular}{|c|c|c|c|c|}
\hline
\rowcolor{gray!15}\textbf{Trung bình} & \textbf{Nhỏ nhất} & \textbf{Lớn nhất} & \textbf{p95} & \textbf{Ngưỡng} \\ \hline
1,858 s & 1,581 s & 3,205 s & 2,155 s & 3,0 s \\ \hline
\end{tabular}
\caption{Tổng hợp 20 lần đo}
\end{table}

\textbf{Kết luận:} \textbf{đạt}, trung bình 1,858 giây thấp hơn ngưỡng 3 giây, và 19/20 lần dưới ngưỡng; chỉ lần đầu tiên vượt ngưỡng (3,205 giây). Thời gian giảm dần theo các lần đầu (khoảng 2,0 giây ở lần 2--8, khoảng 1,6 giây từ lần 10 trở đi); chưa kiểm chứng nguyên nhân, có thể do Moodle/PHP và trình duyệt được làm nóng bộ nhớ đệm. Thời gian mở form chỉ khoảng 0,3--0,8 giây. Ngưỡng 3 giây do nhóm tự đặt vì đề không quy định. 

\textbf{Cấu hình máy đo:} Windows, RAM 16 GB DDR4 (khoảng 15,7 GB khả dụng), đồ họa tích hợp Intel Iris Xe, ổ SSD NVMe. Docker Desktop 4.96.2 dùng WSL 2, Docker \textbf{12 CPU (luồng logic)} và \textbf{7,6 GiB RAM}.

\subsection{Chi tiết phần của }
\subsection{Chi tiết phần của }
\subsection{Chi tiết phần của }
\subsection{Chi tiết phần của Phát: tương thích của Add a Choice}
\textbf{Loại kiểm thử:} kiểm thử tương thích chức năng trên Chrome và Microsoft Edge.
Yêu cầu được lựa chọn là \texttt{NFR-COMP-001}: bốn luồng đại diện của Add a Choice phải chạy đạt
trên cả hai trình duyệt và cho dữ liệu đọc lại tương đương.

\textbf{Cách tiếp cận:} chọn TC-005-001 (tạo hợp lệ), TC-005-002 (thiếu tên rồi sửa),
TC-005-013 (cho phép chọn nhiều phương án) và TC-005-024 (Cancel).
Chạy cùng dữ liệu lần lượt trên Chrome và Edge bằng mã Level 2, sau đó so sánh các quan sát thực tế:
tên hoạt động, cấu hình đã lưu, số phương án, thông báo lỗi hoặc trạng thái danh sách không đổi sau Cancel.
Các ca này đại diện cho bốn luồng, không thay thế bộ 24 ca đã chạy riêng ở mỗi Level.

\textbf{Công cụ:} Python, Selenium WebDriver, \texttt{unittest}, trình duyệt Chrome và Edge.

\begin{lstlisting}[language=Python]
CASES = ["TC-005-001", "TC-005-002", "TC-005-013", "TC-005-024"]
for browser in ("chrome", "edge"):
    subprocess.run([
        sys.executable, "level2/test_choice.py",
        "--browser", browser, "--case", ",".join(CASES),
        "--course-id", course_id
    ])
# Read results.csv and observations.json from each browser run.
# Compare the actual observations for corresponding case IDs.
\end{lstlisting}
Đoạn trên rút gọn cách gọi lại mã kiểm thử. Script đầy đủ nằm ở
\texttt{F005/nonfunctional/test\_compatibility.py}, sử dụng đường dẫn tương đối đến
\texttt{F005/level2/test\_choice.py}; hai thư mục phải được giữ cạnh nhau.

\textbf{Tiêu chí đạt:} cả bốn ca đều pass trên mỗi trình duyệt và các giá trị đọc lại của từng cặp ca
giống nhau. Chỉ một trình duyệt không đạt hoặc quan sát khác nhau thì phần kiểm thử tương thích
không được kết luận pass.

\textbf{Script và cách chạy:} từ thư mục \texttt{F005}, chạy lệnh dưới đây; thay ID 2 bằng ID khóa học
thực tế trên máy đang dùng:
\begin{lstlisting}
python nonfunctional/test_compatibility.py --course-id 2
\end{lstlisting}

\textbf{Kết quả chạy:} Chrome 158.0.8127.37 và Edge 158.0.4370.33,
Selenium 4.54.0, cùng Moodle demo và khóa F005 của lần chạy chức năng.

\begin{table}[H]
\centering\small
\begin{tabular}{|L{2.6cm}|L{5.0cm}|c|c|c|}
\hline
\rowcolor{gray!15}\textbf{Ca} & \textbf{Luồng} & \textbf{Chrome} & \textbf{Edge} & \textbf{So sánh} \\ \hline
TC-005-001 & Tạo Choice hợp lệ & Pass & Pass & Giống nhau \\ \hline
TC-005-002 & Báo lỗi tên rỗng, sửa và lưu & Pass & Pass & Giống nhau \\ \hline
TC-005-013 & Lưu thiết lập chọn nhiều & Pass & Pass & Giống nhau \\ \hline
TC-005-024 & Cancel không thêm hoạt động & Pass & Pass & Giống nhau \\ \hline
\end{tabular}
\caption{Kết quả tương thích Chrome/Edge của Feature 005}
\end{table}

\textbf{Kết luận:} \textbf{đạt} trong phạm vi đã kiểm tra: 8/8 lượt thực thi pass,
4/4 cặp quan sát tương đương; script in \texttt{NFR-COMP-001: PASS}.
File \texttt{compatibility.csv} ghi kết quả so sánh, \texttt{report.json} chứa tổng hợp và các quan sát;
mỗi trình duyệt có thêm \texttt{results.csv} và \texttt{run.json}.
Thư mục bằng chứng trên Drive được đặt tên \texttt{F005\_Compatibility\_Chrome\_Edge}.

Chrome và Edge cùng sử dụng Chromium, nên kết quả chỉ xác nhận bốn luồng nêu trên ở hai phiên bản
đã chạy trên Windows; chưa có kết quả cho Firefox, Safari, thiết bị di động hoặc đánh giá accessibility.


\section{AI Usage Declaration}

\subsection{Nguyễn Đức Phương}
\textbf{Công cụ:} Claude (Anthropic), model Claude Sonnet 5.5, dùng qua ứng dụng Claude (bản desktop).

\noindent\textbf{Prompt đã sử dụng} cho Bài tập lớn 3 (nguyên văn):
\begin{enumerate}[noitemsep]
  \item ``hãy trình bày lại yêu cầu đề cho tôi'' (kèm file mô tả đề)
  \item ``hãy soạn phần report cho tôi''
\end{enumerate}

\noindent\textbf{Kết quả, nhận xét và đề xuất của AI:}
\begin{itemize}[noitemsep]
  \item Tóm tắt yêu cầu đề, thang điểm và tài liệu cần nộp.
  \item Đề xuất cách gộp 50 test case của Feature 001 thành 3 script data-driven.
  \item Soạn khung báo cáo.
\end{itemize}

\subsection{Lê Sĩ Nhật Quang}
\todo{Điền công cụ AI và prompt, hoặc câu xác nhận không dùng AI.}
\subsection{Huỳnh Đức Tài}
\todo{Điền công cụ AI và prompt, hoặc câu xác nhận không dùng AI.}
\subsection{Đào Xuân Sơn}
\todo{Điền công cụ AI và prompt, hoặc câu xác nhận không dùng AI.}
\subsection{Sơn Văn Phát}
\textbf{Công cụ:} GPT-6 (OpenAI), sử dụng qua ứng dụng Codex.

\noindent\textbf{Yêu cầu đã gửi cho AI} (tóm tắt nội dung trao đổi):
\begin{enumerate}[noitemsep]
  \item Trình bày lại mô tả chính của BTL3 và yêu cầu chung về hình thức nộp bài, dựa trên file mô tả đề.
  \item Hỗ trợ chuyển bộ 24 ca Add a Choice sang Python.
\end{enumerate}

\noindent\textbf{Kết quả, nhận xét và đề xuất của AI:}
\begin{itemize}[noitemsep]
  \item Tóm tắt yêu cầu của hai Level, phần kiểm thử phi chức năng và các file cần nộp.
  \item Hỗ trợ viết và rà soát mã Python; đề xuất cách đọc dữ liệu, kiểm tra kết quả đã lưu và dọn hoạt động sau mỗi ca.
  \item Phân tích lỗi thao tác Edit mode, nhập mô tả và chờ xóa dữ liệu; hỗ trợ trình bày báo cáo từ log chạy thực tế.
\end{itemize}

\end{document}
